\scanpage{48}
On a lissité (à tous les coups) ssi \deleted{\unsure{pour}} la suite $i_p\ i_{p+1}\ldots i_n$ (à partir de $i_p$, pas de $i_1$ !) est \unsure{saturée}, i.e. chaque terme \inserted{(sauf le premier)} est successeur \inserted{dans $I$} du précédent.

11. \uline{Digression : où on laisse tomber} $I$.

Soit donné un ens. ordonné topologique, i.e. un \deleted{espace topologique} ens. ordonné $\Sigma=\Sigma(0),\Sigma(1)$, où $\Sigma(0)$ est un espace topologique, $\Sigma(1)$ une partie fermée de $\Sigma(0)\times\Sigma(0)$, qui soit le graphe d'une relation d'ordre. On suppose
\margin{\unsure{On retient ce cas}}

\textcircled{a}\quad $\Sigma(1)\xrightarrow{\mathrm{but}}\Sigma(0)$ est une immersion localement bordante, et propre (donc finie)

\deleted{\textcircled{b}} Cela implique que $\Sigma(0)\xrightarrow{\mathrm{diag}}\Sigma(1)$ est un isomorphisme sur une partie ouverte et fermée. Soit $\Sigma^{<}(1)$ le complémentaire de l'image, où l'on a une des inégalités strictes $i<j$. On désigne par $\Sigma^{<}(n)$ l'espace des drapeaux \uline{stricts} $i_0<i_1\ldots<i_n$ de longueur $n$. Les $\Sigma^{\unsure{*|<}}(n)$ \deleted{\unsure{sont forment|forment}} un espace \inserted{pré}\unsure{semi}-simplicial (\uline{NB} pas d'applications de dégénérescence, \deleted{\unsure{car par}} ils sont dans les $\Sigma(n)\ldots$). \deleted{Les} Si on a un
\mspage{47}
